% --- Содержимое файла materials/5_5.tex ---

\subsection{Поточечная сходимость последовательностей операторов}

\begin{idea}[Смысл теоремы]
	Если последовательность операторов сходится в каждой точке, то предельный оператор «наследует» линейность и непрерывность.
	\\
	\textbf{Важно:} Это работает, только если пространство $E_1$ полное (чтобы работала теорема Банаха-Штейнгауза для оценки норм) и $E_2$ полное (чтобы предел вообще существовал).
\end{idea}

\begin{theorem}[О полноте относительно поточечной сходимости]
	Пусть $E_1, E_2$ --- банаховы пространства.
	Пусть $\{A_n\} \subset \mathcal{L}(E_1, E_2)$ --- последовательность операторов, которая является фундаментальной в каждой точке:
	\[ \forall x \in E_1 \quad \{A_n x\} \text{ --- фундаментальна в } E_2. \]
	Тогда существует предельный линейный непрерывный оператор $A \in \mathcal{L}(E_1, E_2)$ такой, что:
	\[ \forall x \in E_1 \quad \lim_{n \to \infty} A_n x = Ax. \]
\end{theorem}

\begin{proof}
	\textbf{1. Существование и Линейность.}
	Так как для любого $x$ последовательность $\{A_n x\}$ фундаментальна, а пространство $E_2$ полное (банахово), то предел существует. Определим оператор $A$:
	\[ Ax := \lim_{n \to \infty} A_n x. \]
	Линейность $A$ следует из свойств предела и линейности каждого $A_n$.
	
	\textbf{2. Непрерывность (Ограниченность).}
	Так как последовательность $\{A_n x\}$ сходится, она ограничена для каждого $x$.
	Значит, семейство $\{A_n\}$ поточечно ограничено.
	Так как $E_1$ --- банахово, по теореме \textbf{Банаха--Штейнгауза} (Thm \ref{Banach-Steingaus}) оно ограничено равномерно:
	\[ \exists M > 0: \quad \forall n \in \mathbb{N} \quad \|A_n\| \leqslant M. \]
	
	Оценим норму предельного оператора. Для любого $x$:
	\[ \|A_n x\| \leqslant \|A_n\| \cdot \|x\| \leqslant M \|x\|. \]
	Переходя к пределу при $n \to \infty$ (учитывая непрерывность нормы):
	\[ \|Ax\| \leqslant M \|x\|. \]
	Следовательно, $A$ ограничен и $A \in \mathcal{L}(E_1, E_2)$.
\end{proof}

\begin{idea}[Метод плотности]
	Как проверить, что операторы $A_n$ сходятся к $A$ во всём бесконечном пространстве? Проверять каждую точку невозможно.
	Критерий ниже дает алгоритм:
	1. Проверь \textbf{стабильность}: нормы операторов не должны расти бесконечно ($\|A_n\| \le M$).
	2. Проверь \textbf{сходимость на простых векторах} (на базисе, многочленах и т.д.).
	Если оба условия выполнены — сходимость есть везде.
\end{idea}


\begin{theorem}[Критерий поточечной сходимости]\label{bshetengcorollary}
	Пусть $E_1$ --- банахово, $E_2$ --- линейное нормированное пространство.
	Пусть $\{A_n\} \subset \mathcal{L}(E_1, E_2)$ и $A \in \mathcal{L}(E_1, E_2)$.
	
	Последовательность $A_n$ сходится к $A$ поточечно ($\forall x: A_n x \to Ax$) тогда и только тогда, когда выполнены два условия:
	\begin{enumerate}
		\item Последовательность норм ограничена: $\exists M: \sup_n \|A_n\| \leqslant M$.
		\item Сходимость есть на плотном множестве: $\exists S \subset E_1$, $\overline{[S]} = E_1$ такое, что $\forall s \in S: A_n s \to As$.
	\end{enumerate}
\end{theorem}

\begin{proof}
	\textbf{Необходимость ($\Rightarrow$)}
	\begin{itemize}
		\item Сходимость на $S$ следует из сходимости на всем $E_1$ (так как $S \subset E_1$).
		\item Ограниченность норм следует из теоремы Банаха--Штейнгауза (так как есть поточечная сходимость на всем $E_1$, значит есть и поточечная ограниченность).
	\end{itemize}
	
	\textbf{Достаточность ($\Leftarrow$)}
	Пусть условия 1 и 2 выполнены. Зафиксируем произвольный $x \in E_1$ и $\varepsilon > 0$.
	Так как линейная оболочка $[S]$ плотна в $E_1$, найдем элемент $y \in [S]$, близкий к $x$:
	\[ \|x - y\| < \frac{\varepsilon}{3M^*}, \quad \text{где } M^* = \max(M, \|A\|). \]
	
	Оценим разность $\|A_n x - Ax\|$ методом «$\varepsilon/3$»:
	\[ \|A_n x - Ax\| \leqslant \underbrace{\|A_n x - A_n y\|}_{\text{(1)}} + \underbrace{\|A_n y - Ay\|}_{\text{(2)}} + \underbrace{\|Ay - Ax\|}_{\text{(3)}}. \]
	
	Разберем слагаемые:
	\begin{itemize}
		\item (1) $\|A_n(x - y)\| \leqslant \|A_n\| \|x - y\| \leqslant M \cdot \frac{\varepsilon}{3M^*} \leqslant \frac{\varepsilon}{3}$.
		\item (3) $\|A(y - x)\| \leqslant \|A\| \|x - y\| \leqslant \|A\| \cdot \frac{\varepsilon}{3M^*} \leqslant \frac{\varepsilon}{3}$.
		\item (2) Так как $y \in [S]$, то $A_n y \to Ay$. Значит, $\exists N$, что при $n > N$: $\|A_n y - Ay\| < \frac{\varepsilon}{3}$.
	\end{itemize}
	
	Суммируя, получаем:
	\[ \|A_n x - Ax\| < \frac{\varepsilon}{3} + \frac{\varepsilon}{3} + \frac{\varepsilon}{3} = \varepsilon. \]
	Значит, $A_n x \to Ax$ для любого $x$.
\end{proof}

\begin{idea}
	Этот критерий лежит в основе обоснования многих численных методов (например, сходимости квадратурных формул). Мы проверяем сходимость для полиномов (это легко), убеждаемся в ограниченности норм сумм, и получаем сходимость для любой непрерывной функции.
\end{idea}